% !TeX spellcheck = ru_RU
% !TEX root = vkr.tex

\section{Постановка задачи}
\label{sec:task}

Конечным пользователем решения является пользователь командной строки,
которому необходим небольшой архиватор и воспроизводимая реализация алгоритма
Хаффмана для учебных экспериментов. Пользователь должен иметь возможность
сжать произвольный файл, восстановить его без потерь и повторить измерения на
опубликованном коде.

\textbf{Целью работы является} реализация архиватора файлов на основе
статического кодирования Хаффмана и экспериментальная проверка его свойств.
Для достижения цели были поставлены следующие задачи:
\begin{enumerate}
    \item выполнить обзор алгоритма Хаффмана и его применения в существующих
          средствах сжатия;
    \item спроектировать формат архива и реализовать сжатие и распаковку
          текстовых и бинарных файлов;
    \item обеспечить проверку корректности реализации модульными и
          интеграционными тестами, а также средствами автоматического анализа;
    \item провести измерения коэффициента сжатия и времени обработки на данных
          с различными свойствами и проанализировать результаты.
\end{enumerate}

Результатом работы должен быть опубликованный репозиторий консольной утилиты
с инструкцией по сборке и запуску, описанием формата архива, автоматическими
проверками и воспроизводимым сценарием измерений.
